% II. Анализ на възможните решения и обосновка на избора.
% Всеки слой има повече от един разумен вариант. Тук се излагат вариантите,
% критериите и изборът. Числените сравнения идват след измерване, не оттук.
\section{Анализ на възможните решения и обосновка на избора}
\label{sec:alternatives}

\subsection{Критерии за избор}

Решенията се оценяват по пет критерия: покритие на конструкциите от предметната
област, поведение при нееднозначност, обяснимост на резултата, цена на пренасяне към
нова предметна област и сложност на реализацията в рамките на курсовата работа.
Обяснимостта се извежда като отделен критерий, защото системата връща на потребителя
перифраза на разбрания въпрос, а такава перифраза е възможна само ако тълкуването е
достъпно като структура, а не само като краен низ.

\subsection{Морфологичен слой: преобразувател срещу речник от словоформи}

Първият вариант е пълен списък от словоформи с приписани признаци. Той е тривиален за
реализация и бърз при търсене, но расте линейно с лексикона и не обработва дума, която
не е в списъка. Вторият вариант е краен преобразувател над лексикон и правила за
словоизменение~\cite{beesley2003fsmorphology}. Той изисква повече работа в началото,
но е компактен, работи в двете посоки и позволява анализ на непозната дума по нейното
окончание. Изборът е вторият вариант, защото синтезът на изречение в
Раздел~\ref{sec:design} се нуждае точно от обратната посока на същото устройство.

\subsection{Синтактичен слой: Earley срещу CYK}

Двата разгледани алгоритъма за анализ при нееднозначност са
CYK~\cite{younger1967cyk} и Earley~\cite{earley1970parsing}. CYK е по-прост за
реализация и има предсказуемо поведение, но изисква граматиката да бъде приведена в
нормална форма на Чомски, което изкривява дървото на извода и усложнява семантичното
изображение. Earley работи с произволна безконтекстна граматика, обработва изречението
инкрементално и дава същата асимптотична граница в общия случай, като е по-бърз при
еднозначни фрагменти.

\TODO{окончателен избор между Earley и CYK след измерване върху набора от въпроси} При
всички положения изходът на слоя е компактна гора от изводи, а не едно дърво, защото
изборът между разчитанията се прави едва след семантичната проверка.

\subsection{Модел за представяне на знания}

Съпоставката между релационен модел, фреймово представяне и графов модел от типа RDF
се прави по критериите от подраздел~2.1. Релационният модел е близък до целевата
заявка на SQL и не изисква допълнително изображение, но описва слабо йерархията от
понятия. Графовият модел описва добре йерархията и се запитва със SPARQL, но налага
допълнителен слой между семантиката и данните. Оформената съпоставка е дадена в
\TODO{таблица със съпоставителна оценка на моделите за представяне на знания}, а
проектът поддържа и двете цели, за да може разликата да бъде измерена, а не постулирана.

\subsection{Символно решение срещу базово решение с голям езиков модел}

Голям езиков модел произвежда SQL от въпрос без нито един от изброените слоеве. Тази
алтернатива е разгледана изрично, защото е очевидната съвременна съперница на
символния подход. Тя има по-широко покритие и не изисква лексикон, но има три
съществени недостатъка за настоящата задача: не дава проверима следа от входа до
изхода, не може да обоснове защо е избрала едно разчитане пред друго, и се проваля по
начин, който потребителят не може да забележи, тъй като изходът изглежда правдоподобен
и когато е грешен.

Точно този контраст прави съпоставката интересна. Символната система има по-тясно
покритие, но всяко нейно тълкуване е представимо, обяснимо и потвърдимо от потребителя
преди изпълнение. Затова езиковият модел се използва тук като \term{базово решение} за
сравнение върху един и същ набор от въпроси, а не като част от системата. Начинът, по
който съпоставката се провежда, е описан в Раздел~\ref{sec:method}.
